\section{Two Sacrifice Channels: Money and Time}
\label{sec:two_channels}

The sacrifice dimension splits into two observation channels.
Money-sacrifice and time-sacrifice are not disjoint coverage
spaces; they are \emph{different modalities of the same
underlying sacrifice substrate}, connected by cooperative-market
exchange.  The value of observing both is triangulation and
robustness, with genuine strict inequality arising only at
specific non-fungibility residuals.

\subsection{Money-sacrifice channel}

\begin{definition}[Money-Sacrifice Event]
\label{def:money_sacrifice}
A money-sacrifice event is a revealed-sacrifice event
$(i, X, Y, t)$ where the benchmarked sacrifice $X$ is purchasing
power: agent~$i$ pays a price $p$ for an unbenchmarked bundle $Y$.
The signal is:
\[
  \volRlower(Y) \;\geq\; \Delta\volL(p),
\]
where $\Delta\volL(p)$ is the $\volL$ contribution of the
purchasing power surrendered.
\end{definition}

\noindent\textbf{Properties of the money channel:}
\begin{itemize}
\item \textbf{Signal strength.}  Price is a precise, quantitative
  signal.  Market aggregation handles distributional loading (many
  agents trading the same category of goods produces a robust
  estimate).
\item \textbf{Bias.}  Privileges capabilities valued by wealthier
  agents.  An agent with more purchasing power emits stronger
  signals per trade.
\item \textbf{Third-party observability.}  Money-sacrifice events
  leave public ledger traces (prices, timestamps, category
  labels); auditors can compute the money-channel aggregate
  without agent cooperation (Property~P5 of
  Proposition~\ref{prop:privacy_institutional}).
\end{itemize}

\subsection{Time-sacrifice channel}

\begin{definition}[Time-Sacrifice Event]
\label{def:time_sacrifice}
A time-sacrifice event is a revealed-sacrifice event
$(i, X, Y, t)$ where the benchmarked sacrifice $X$ is the agent's
outside-option labour hours, valued at market wage rate $r_i$.
Agent~$i$ foregoes $r_i \cdot h$ hours of wage income to engage in
an unbenchmarked activity (volunteering, parenting, craft practice,
relationship-building, contemplation).  The signal is:
\[
  \volRlower(Y) \;\geq\; \Delta\volL(r_i \cdot h).
\]
\end{definition}

\noindent\textbf{Properties of the time channel:}
\begin{itemize}
\item \textbf{Signal strength.}  Opportunity cost is a
  quantitative signal grounded in market wage rates.  Every agent
  has a time endowment and an outside option, regardless of
  wealth.
\item \textbf{Bias.}  More egalitarian than the money channel.
  Time is the scarcest capability every agent shares.  A
  high-wage agent's time sacrifice has higher $\Delta\volL$ per
  hour, but every agent sacrifices time.
\item \textbf{Non-market agents.}  Agents not in the labour
  market (retirees, children, homemakers, the unemployed) have
  an undefined or zero market wage rate.  The framework can use
  imputed wage rates~\citep{aguiar2007measuring} or non-market
  time valuation (replacement cost) to extend coverage to these
  populations.
\item \textbf{Public versus private sacrifice.}  Publicly
  observable time-sacrifice (committed volunteer hours,
  time-stamped artefact production) is third-party observable.
  Private time-sacrifice (parenting, contemplation) is visible
  only via voluntary disclosure---the residual flagged in
  Property~P5 of Proposition~\ref{prop:privacy_institutional}.
\end{itemize}

\subsection{Substitutability via cooperative markets}
\label{sec:channel_substitutability}

Money and time are not disjoint sacrifice types.  Cooperative
markets exchange between them:
\begin{itemize}
\item \textbf{Time $\to$ money.}  Employment converts labour hours
  to wage income at rate $r_i$.  This is the definitional
  exchange rate already embedded in
  Definition~\ref{def:time_sacrifice}.
\item \textbf{Money $\to$ time.}  Agents can hire others to
  invest labour hours (services, delegation) or purchase
  time-saving technology and infrastructure (automation,
  efficiency-enhancing tools).  The exchange goes through the
  same cooperative market that sustains employment.
\end{itemize}

Because the two channels are linked by this exchange, an agent's
total sacrifice budget can be measured in a common monetary unit
via the wage rate, and the agent can shift sacrifice between
modalities.  A high-earning agent can substitute money for time
by hiring help; a time-rich agent can substitute time for money
by investing hours directly.  Neither channel \emph{dominates}
the other---each is the natural view of the same underlying
sacrifice from one side of the exchange.

The practical consequence: \textbf{the two channels triangulate
the same capability space, they do not partition disjoint
capability spaces}.  Observing both yields:

\begin{itemize}
\item \textbf{Triangulation.}  When money-sacrifice and
  time-sacrifice both indicate high $\volRlower$ on a bundle,
  the signals corroborate.  Divergence between the channels is
  itself diagnostic: a capability with strong time-sacrifice
  evidence but no money-sacrifice evidence in a given window
  may indicate the capability is exercised by time-rich agents
  outside the monetary economy, or that market mechanisms have
  not formed around it.
\item \textbf{Robustness.}  Different agent populations emit
  different modal signals.  Children, retirees, and non-market
  agents emit time signals with imputed wage rates; wage-
  constrained high earners emit money signals.  Observing both
  modalities extends the bound across populations that either
  alone would miss.
\end{itemize}

\subsection{Non-fungibility residuals}

The substitutability in \S\ref{sec:channel_substitutability} is
not total.  Two specific residuals create genuine asymmetries
where one channel captures something the other cannot substitute
for.

\begin{proposition}[Agent-Specific-Time Residual (existential)]
\label{prop:agent_specific_time}
There exists a capability $c^\ast$ reachable by the time channel
but whose $\volR$-content is not captured by any money-channel
event on $c^\ast$ itself (access goods purchasable with money
attach $\volR$-content to \emph{access categories}, not to
$c^\ast$).  The non-substitutability is of the specific
$c^\ast$, not a claim that all agent-specific-time capabilities
evade the money channel entirely.
\end{proposition}

\begin{proof}
We exhibit $c^\ast = $ ``sustained mentoring relationship with a
specific individual.''  The relational capability is constituted
by the interaction between this agent and the mentee; it is not
transferable to a proxy.  Money can purchase access goods
(introductions, a mentor's scheduled time, coaching
infrastructure), but these are distinct bundles whose
$\volR$-content lies in access categories rather than in the
relational capability $c^\ast$ itself.  A paid mentoring
engagement generates a money-channel event on \emph{scheduled
access} rather than on the relationship capability $c^\ast$.
When agent~$i$ allocates $h$ hours foregoing wage income
$r_i \cdot h$ to the relationship, the time-sacrifice event
(Definition~\ref{def:time_sacrifice}) produces
$\volRlower(Y) \geq \Delta\volL(r_i \cdot h)$ with $c^\ast \in
Y$, giving the time channel an event that the money channel
does not produce for $c^\ast$.
\end{proof}

\begin{proposition}[Capital-Concentrated Residual (existential)]
\label{prop:capital_concentrated}
There exists a capability $c^\ast$ reachable by the money
channel but whose $\volR$-content is not captured by any
time-channel event on $c^\ast$ itself (direct time-sacrifice on
$c^\ast$ without capital access yields no realized exercise).
The non-substitutability is of the specific $c^\ast$, not a
claim that all capital-concentrated capabilities evade the time
channel entirely.
\end{proposition}

\begin{proof}
We exhibit $c^\ast = $ ``conduct a particle-physics experiment
using a synchrotron.''  The exercise of $c^\ast$ requires access
to specialized infrastructure whose acquisition or lease cost is
concentrated in a single transaction (institutional budget,
grant award, facility time purchased).  An agent accumulating
time-sacrifice at wage rate $r_i$ can only reach this capability
by first converting time to money via employment and then
purchasing access---which is the money channel, not the time
channel, since the converted-to-money event is a market
transaction.  Direct time-sacrifice (hours spent without capital
access) yields no experimental event and no $\volR$ signal on
$c^\ast$.  The money channel captures $c^\ast$ via the single
access-purchase event; the time channel is silent for $c^\ast$
because the time-sacrifice exchange rate on direct-labor hours
is too small relative to the capital threshold, and
time-channel events accompanying capital-enabled exercise are
attributable to labor categories rather than to $c^\ast$'s
capital-concentrated access.
\end{proof}

\begin{remark}[Scope of the residual propositions]
\label{rem:residual_scope}
Propositions~\ref{prop:agent_specific_time}
and~\ref{prop:capital_concentrated} are
\emph{existential}---each witnesses one capability exhibiting
the claimed residual behavior.  They do not establish that the
descriptive class ``all agent-specific-time capabilities'' or
``all capital-concentrated capabilities'' is uniformly
channel-exclusive.  Mixed-mode coverage of such capabilities can
and does occur (a paid mentor also incurs time-sacrifice; a
well-funded scientist can generate labor-attributed time events
alongside capital-access events).  The existential residuals
are sufficient for Corollary~\ref{cor:joint_coverage}'s
strict-improvement claim, which is stated conditional on the
trade window actually containing at least one residual-category
event.
\end{remark}

\begin{corollary}[Joint Channel Triangulation]
\label{cor:joint_coverage}
Let $E_{\mathrm{money}}$ and $E_{\mathrm{time}}$ denote the
money-channel and time-channel event sets over a trade window
$W$, and suppose the union $E = E_{\mathrm{money}} \cup
E_{\mathrm{time}}$ satisfies the closure conditions of
Theorem~\ref{thm:aggregate_bound} Part~A: the bundles
$\{Y_n\}_{n \in E}$ are pairwise poset-independent across
channels---no subsumption relation and no cooperative-composition
link connects any capability in a money-channel bundle to any
capability in a time-channel bundle (set-disjointness of
bundles across channels is necessary but not sufficient).

\medskip\noindent
\textbf{Weak containment (unconditional).}
The joint aggregate captures \emph{weakly more} of $U$ than
either channel alone: any event observed through either channel
alone is in the union, and Part~A aggregation is monotone in
observed admissible events (Proposition~\ref{prop:monotone}).

\medskip\noindent
\textbf{Strict improvement (conditional on observed coverage).}
The non-fungibility residuals
(Propositions~\ref{prop:agent_specific_time}
and~\ref{prop:capital_concentrated}) make strict improvement
\emph{possible}: agent-specific-time capabilities lie in the
coverage gap of the money channel alone, and
capital-concentrated capabilities lie in the coverage gap of the
time channel alone.  Strictness over the money channel alone
obtains when $W$ contains at least one admissible time-only
residual event; strictness over the time channel alone obtains
when $W$ contains at least one admissible money-only residual
event.  The residuals establish the \emph{availability} of
strict improvement, not its automatic realization; whether any
given window realizes strictness is a data-coverage question.
\end{corollary}

\begin{proof}[Proof sketch]
Weak containment under the stated closure hypothesis: any event
observed through either channel alone is also observed through
the union, and the cross-channel pairwise poset-independence
ensures Part~A's aggregation is well-defined on $E$; the Part~A
lower bound is monotone in observed admissible events under
sub-batch supremum
(Proposition~\ref{prop:monotone_part_a}), so the union's bound
is $\geq$ either single channel's bound.  Strict improvement on the
residuals: Proposition~\ref{prop:agent_specific_time} witnesses
a capability captured by time but not by money;
Proposition~\ref{prop:capital_concentrated} witnesses a
capability captured by money but not by time.  Both residuals
guarantee strict improvement of the joint aggregate over either
channel alone when the trade window includes at least one event
from the relevant residual category.  Outside the residuals,
strictness becomes a data-coverage statement: the joint
aggregate strictly exceeds either single channel if and only if
the trade window includes events in categories the other
channel does not cover.  Substitutability in principle does not
guarantee equality in practice; heterogeneous agent populations
generically produce strict improvement, but a theoretical
population using one channel exclusively could produce
single-channel coverage equal to the joint aggregate.
\end{proof}

\begin{remark}[Cross-channel poset coupling]
\label{rem:cross_channel_overlap}
The closure hypothesis rules out all three forms of
cross-channel coupling: (a)~a capability $c$ appearing in both
a money-channel bundle and a time-channel bundle (set-overlap),
(b)~subsumption across channels ($c \succeq c'$ with $c$ in one
channel's bundle and $c'$ in the other's), and (c)~cooperative
composition linking capabilities across channels.  When any of
these occurs in practice, the operational responses of
Remark~\ref{rem:label_collision} apply: partition events into
maximal poset-independent sub-batches per channel and aggregate
separately, or fall back to Part~B's per-capability
max-attribution (which natively handles overlapping events via
the max across events covering each capability, subject to
Part~B's additional hypotheses).
\end{remark}

\subsection{The residual intersection}

What neither channel catches is the \emph{residual class}
(characterized in the companion paper~\citep{lasser2026paper8b}):
capabilities exercised purely privately, with no trade event,
no publicly-observable time expenditure, and no sacrifice
footprint of any kind.  The
residual class is the intersection of the coverage gaps of both
channels after accounting for the cooperative-market
substitutability of \S\ref{sec:channel_substitutability} and the
non-fungibility residuals of
Propositions~\ref{prop:agent_specific_time}
and~\ref{prop:capital_concentrated}.  This is a strictly smaller
set than the union of the two channels' individual coverage gaps:
capabilities in one channel's gap but reachable via the other
(or via the substitution markets) are not residual.
